% Week naming rule:
% 2026-05-w04 means the Monday of this report week is in May 2026,
% and it is the 4th Monday in that month. This week starts on 2026-05-25.

\begin{frame}
  \titlepage
  \vspace{-1.0em}
  \begin{center}
    \small Reporting period: May 2026, Week 4 \\
    \scriptsize Month/week is determined by the Monday of the reporting week.
  \end{center}
\end{frame}

\begin{frame}{One-Slide Summary}
  \begin{itemize}
    \item We are narrowing Plan 08 into a practical v0: a frozen 7B/8B base model plus a trainable memory wrapper.
    \item The wrapper compresses long context into an updateable state:
      \[
        m_t = G_\phi(m_{t-1}, c_t), \qquad y_t = F_\theta(q_t, m_t)
      \]
    \item We do not directly fine-tune the base model in v0, because we do not have the original pretraining mixture and want to avoid catastrophic forgetting.
    \item Two deliverables:
      \begin{enumerate}
        \item Research paper: new architecture for next-state memory prediction.
        \item Open-source RCA-Code model: public debugging / trace-based RCA model with wrapper memory.
      \end{enumerate}
  \end{itemize}
\end{frame}

\begin{frame}{Motivation: Long Context Is Not Enough}
  \begin{itemize}
    \item RCA and debugging tasks often contain long evidence bundles:
      logs, traces, stack traces, test failures, code context, topology, and prior observations.
    \item Direct long-context prompting has three problems:
      \begin{itemize}
        \item high prefill and KV-cache cost;
        \item lost-in-the-middle / context rot;
        \item no persistent learning between chunks or turns.
      \end{itemize}
    \item Direct SFT on a small RCA/debug dataset risks degrading the base model.
    \item Our hypothesis: learn a small wrapper memory while keeping the base model frozen.
  \end{itemize}
\end{frame}

\begin{frame}{Core Architecture}
  \begin{columns}[T]
    \begin{column}{0.52\textwidth}
      \begin{block}{Memory update}
        \[
          m_t = G_\phi(m_{t-1}, c_t)
        \]
        \[
          y_t = F_\theta(q_t, m_t)
        \]
      \end{block}
      \begin{itemize}
        \item $F_\theta$: frozen 7B/8B base model.
        \item $G_\phi$: trainable wrapper.
        \item $m_t$: compact memory state.
        \item $c_t$: new context chunk or tool observation.
      \end{itemize}
    \end{column}
    \begin{column}{0.44\textwidth}
      \begin{block}{V0 memory type}
        Start with explicit memory tokens.
      \end{block}
      \begin{itemize}
        \item Easy to train and ablate.
        \item Compatible with frozen decoder models.
        \item Safer than writing into weights.
        \item Hidden-state memory and LoRA memory are later ablations.
      \end{itemize}
    \end{column}
  \end{columns}
\end{frame}

\begin{frame}{Why Freeze the 8B Base?}
  \begin{itemize}
    \item We start around 8B because it is strong enough for reasoning and still feasible for repeated experiments.
    \item Candidate bases:
      \begin{itemize}
        \item Qwen3-8B / Qwen3-8B-Instruct for long-context and reasoning experiments.
        \item Llama-3.1-8B-Instruct as a secondary validation base.
        \item Qwen2.5-Coder-7B or similar for RCA-Code.
      \end{itemize}
    \item The base remains frozen in v0.
    \item Training only the wrapper limits catastrophic forgetting and makes the contribution easier to isolate.
  \end{itemize}
\end{frame}

\begin{frame}{Two Tracks}
  \begin{columns}[T]
    \begin{column}{0.48\textwidth}
      \begin{block}{Track A: Research Paper}
        \begin{itemize}
          \item New memory-wrapper architecture.
          \item Next-state prediction over memory states.
          \item Evaluate on long-context memory benchmarks.
          \item Main claim: compact learned memory can replace part of long prompt context.
        \end{itemize}
      \end{block}
    \end{column}
    \begin{column}{0.48\textwidth}
      \begin{block}{Track B: RCA-Code Model}
        \begin{itemize}
          \item Public debugging / trace RCA model.
          \item Use code issues, stack traces, failing tests, patches.
          \item Direct prediction and agentic tool-use modes.
          \item Release model, report, and demo.
        \end{itemize}
      \end{block}
    \end{column}
  \end{columns}
\end{frame}

\begin{frame}{Dataset Plan}
  \begin{table}
    \small
    \begin{tabular}{p{0.26\textwidth} p{0.28\textwidth} p{0.36\textwidth}}
      \toprule
      Track & Datasets & Purpose \\
      \midrule
      Base long-context & LoCoMo, LongBench, RULER, Needle-in-a-Haystack, Qasper, NarrativeQA & Prove wrapper memory before domain-specific RCA. \\
      RCA-Code & SWE-bench Verified, BugsInPy, Defects4J, QuixBugs, curated CodeNet/APPS traces & Learn debugging and trace-based root-cause reasoning from public data. \\
      RCA domain & Nezha, OpenRCA-500, RCAEval, lincyaw/rca & Transfer compression method to RCA evidence bundles. \\
      Internal probe & Liang / Nokia DTS & Qualitative validation only; not for public release. \\
      \bottomrule
    \end{tabular}
  \end{table}
\end{frame}

\begin{frame}{Evaluation Modes}
  \begin{block}{Direct prediction}
    Input: issue text, code context, stack trace, logs, failing test output.\\
    Output: root cause, evidence, fix recommendation.
  \end{block}

  \begin{block}{Agentic RCA}
    The model can call tools such as search, open file, run tests, and inspect traces.\\
    Test-time training updates the wrapper from tool observations; the frozen base does not change.
  \end{block}

  \begin{block}{Baseline comparison}
    Full context vs. summary vs. retrieval vs. learned memory.
  \end{block}
\end{frame}

\begin{frame}{Success Metrics}
  \begin{columns}[T]
    \begin{column}{0.48\textwidth}
      \begin{block}{Quality}
        \begin{itemize}
          \item RCA / debug root-cause accuracy.
          \item Fix recommendation quality.
          \item Evidence grounding.
          \item JSON / schema stability for RCA outputs.
        \end{itemize}
      \end{block}
    \end{column}
    \begin{column}{0.48\textwidth}
      \begin{block}{Compression}
        \begin{itemize}
          \item Token and KV-cache reduction.
          \item Latency reduction.
          \item Early / middle / late evidence retention.
          \item Robustness to chunk order and distractors.
        \end{itemize}
      \end{block}
    \end{column}
  \end{columns}
\end{frame}

\begin{frame}{Budget Estimate}
  \begin{table}
    \small
    \begin{tabular}{p{0.25\textwidth} r r r p{0.24\textwidth}}
      \toprule
      Scope & Time & GPUh & Cost & Output \\
      \midrule
      Lean MVP & 4--7 wks & 1,420--2,950 & \$4.3k--\$8.9k & One wrapper, smoke + RCA-Code demo. \\
      Paper-quality v0 & 7--13 wks & 2,820--5,750 & \$8.5k--\$17.3k & Long-context + RCA-Code + ablations. \\
      Strong release & 10--16 wks & 5,000--9,000 & \$15k--\$27k & Multiple seeds, model card, polished report. \\
      \bottomrule
    \end{tabular}
  \end{table}

  \vspace{0.4em}
  Cost model: H100 at \$3 per GPU-hour.
\end{frame}

\begin{frame}{Execution Timeline}
  \begin{table}
    \small
    \begin{tabular}{p{0.19\textwidth} p{0.28\textwidth} p{0.42\textwidth}}
      \toprule
      Phase & Duration & Decision gate \\
      \midrule
      0. Smoke & 3--5 days & Wrapper beats same-length summary on synthetic / NIAH. \\
      1. Long-context & 2--3 weeks & Beats summary on LoCoMo / LongBench / RULER. \\
      2. RCA-Code & 2--4 weeks & Improves direct root-cause / fix prediction on public debug traces. \\
      3. RCA domain & 1--2 weeks & Works on public RCA datasets with at least 4x compression. \\
      4. Release polish & 1--2 weeks & Ablations, demo, model card, technical report. \\
      \bottomrule
    \end{tabular}
  \end{table}
\end{frame}

\begin{frame}{Current Status}
  \begin{itemize}
    \item Plan 08 v0 has been scoped as a learned-memory wrapper project.
    \item English and Chinese notes are maintained in the research notes repo.
    \item Budget is separated from the full self-modifying-LLM version.
    \item RCA-Code path is defined around public coding/debug datasets.
    \item Initial RCA compression scripts and report templates are drafted in the RCA demo repo.
  \end{itemize}
\end{frame}

\begin{frame}{Next Steps}
  \begin{enumerate}
    \item Build the first LoCoMo / NIAH smoke benchmark for wrapper memory.
    \item Run frozen 8B base baselines: full context, summary, retrieval.
    \item Train the first explicit memory-token wrapper.
    \item Convert public coding/debug data into RCA-Code format.
    \item Prepare a small demo comparing full context vs. learned memory.
  \end{enumerate}
\end{frame}

\begin{frame}{Ask / Discussion}
  \begin{itemize}
    \item Is frozen-base + wrapper the right first milestone?
    \item Which 8B base should be the first public model target?
    \item Should the first paper emphasize long-context memory, RCA-Code, or both?
    \item Which demo case would be most convincing for external readers?
  \end{itemize}
\end{frame}
